\section{Time derivatives and variations of whole trajectories}
\label{sec:trajectory_variations}
This discussion distinguishes two sources of differential observations:
time evolution along one trajectory and variations across a family of
admissible trajectories. Independent perturbation trials are not a general
requirement for online learning. Their role is to supply additional
directions when temporal observations do not identify the desired partials.
The calculations below are explanatory examples, not new simulation results
or evidence that the retained structured-coefficient studies learn an
unrestricted differential field.

\subsection{A trajectory variation is not an arbitrary coordinate change}
Let $f:\mathbb R^d\to\mathbb R^m$ be continuously differentiable, with
$F(z)=D_zf(z)$ of row rank $m<d$ on the region of interest. All explicit
time or phase dependencies are included among the context coordinates.
Consider a smooth family of complete admissible trajectories
$z_\varepsilon(t)$ satisfying $f(z_\varepsilon(t))=0$, with
$z_0(t)=z(t)$. Differentiation with respect to the trial parameter gives
\begin{equation}
 \delta z(t)=\left.\frac{\partial z_\varepsilon(t)}{\partial\varepsilon}
                    \right|_{\varepsilon=0},
 \qquad F(z(t))\delta z(t)=0.
 \label{eq:tv_family}
\end{equation}
The parameter changes a whole input history and, where permitted, initial
conditions; the physical system determines the corresponding state and
output histories. Context channels need not be independently commandable.
For example, if $r_\varepsilon=\partial_t x_\varepsilon$, smoothness gives
$\delta r=\partial_t\delta x$. A fixed clock or fixed initial condition
likewise constrains the admissible variations. Independent directions here
mean linearly informative admissible variations, not arbitrary independent
changes to every component of $z$ or probabilistic independence of trials.

Differentiation with respect to physical time instead gives
\begin{equation}
 v(t)=\dot z(t),\qquad F(z(t))v(t)=0.
 \label{eq:tv_temporal}
\end{equation}
Locally, the shifted family $z_\varepsilon(t)=z(t+\varepsilon)$ has
$\delta z=\dot z$. It need not be an admissible comparison trial under fixed
initial conditions, endpoints, or clock phase; the temporal chain-rule
identity in (\ref{eq:tv_temporal}) still holds. If explicit time dependence
is omitted from $z$, the corresponding $f_t$ term must instead be retained.
Physical time therefore provides one tangent direction at each observed
context, whereas admissible trajectory families can provide other directions.

Both observations can enter the direct incompatibility $\xi=Qv$ with all
entries of $Q$ adaptive. The observation type, context units, parameter
normalization, and memory weights must be stated; a trial derivative and a
time derivative should not be pooled with unspecified relative scaling.
True partials are not learning targets supplied to this update.

\subsection{When temporal history is sufficient}
One direction at one instant does not determine a general differential.
Several times can nevertheless identify a shared local row space. To see
the exact requirement, suppose a full-row-rank matrix $F_0$ is compatible
with every retained observation $v_i$, and form the uncentered second moment
\begin{equation}
 F_0v_i=0,\qquad
 R=\sum_i\alpha_i v_iv_i^{\mathsf T},\quad \alpha_i>0.
 \label{eq:tv_covariance}
\end{equation}
Here $R$ has no additional ridge or incompatible initial-memory term.
For any column vector $q$,
$q^{\mathsf T}Rq=\sum_i\alpha_i(q^{\mathsf T}v_i)^2$.
Consequently $\ker R$ is the orthogonal complement of the observed span.
Because that span lies in $\ker F_0$, full tangent excitation gives
\begin{equation}
 \operatorname{span}\{v_i\}=\ker F_0
 \quad\Longleftrightarrow\quad \operatorname{rank}R=d-m.
 \label{eq:tv_identification}
\end{equation}
Under these equivalent conditions,
$\ker R=\operatorname{range}(F_0^{\mathsf T})$.
This identifies the row space; a basis remains ambiguous under nonsingular
row transformations. It is an identifiability statement, not a guarantee
that every initialization or adaptive update converges to that space.

For a changing differential, different times need not share one exactly
compatible $F_0$. Online local tracking therefore requires informative
directions to accumulate on a timescale over which the target row space
changes sufficiently little relative to the excitation gap and adaptation
rate. Finite memory, derivative estimation, and sensor noise create distinct
bias and tracking errors. An empirical spectral gap alone does not establish
that its bottom eigenspace is the physical row space. An appropriate shared
lift can represent faster variation, but its own compatibility, excitation,
and evaluated-field nondegeneracy conditions must hold. An arbitrary
context-dependent field is not identified merely by observing one tangent
at each new context.

\subsection{A missing direction that can be supplied either way}
For illustration only, let a simulator obey
\begin{equation}
 f(x,u,y)=y-2x-3u=0,\qquad z=(x,u,y)^{\mathsf T}.
 \label{eq:tv_linear_example}
\end{equation}
Neither this formula nor any of its partial entries is supplied to the
learner. Along $x(t)=u(t)=t$, the output is $y(t)=5t$, so
\begin{equation}
 \dot z=(1,1,5)^{\mathsf T},\qquad
 F_\star=(-2,-3,1),\quad \widetilde F=(-4,-1,1).
 \label{eq:tv_ambiguous_rows}
\end{equation}
Both rows annihilate every observed time derivative, although they disagree
on the responses to separate changes of $x$ and $u$. The final entry one is
only a convenient convention for writing these two candidates; it is not
fixed in a direct learner. Normalizing either row preserves its zero
incompatibility and does not resolve the ambiguity.

A smooth perturbation of the entire commanded $x$ trajectory gives
\begin{equation}
 \begin{aligned}
 x_\varepsilon(t)&=t+\varepsilon h(t),&u_\varepsilon(t)&=t,\\
 y_\varepsilon(t)&=5t+2\varepsilon h(t),&
 \delta z(t)&=h(t)(1,0,2)^{\mathsf T}.
 \end{aligned}
 \label{eq:tv_linear_trial}
\end{equation}
Choose $h$ to satisfy any imposed endpoint conditions. Where $h(t)\ne0$,
$F_\star\delta z=0$ but $\widetilde F\delta z=-2h(t)\ne0$.
This variation and the temporal direction in
(\ref{eq:tv_ambiguous_rows}) span the two-dimensional physical tangent
plane, identifying its normal row up to scale.

Separate trials are unnecessary for this constant-field example if a single
online input trajectory already excites both directions. Taking $x=t$ and
$u=t^2$ gives
\begin{equation}
 \dot z(t)=(1,2t,2+6t)^{\mathsf T}.
 \label{eq:tv_online_excitation}
\end{equation}
At any two distinct times these vectors are linearly independent. They
annihilate the same constant $F_\star$, so temporal observations alone
satisfy the exact span requirement. This construction explains why extra
trials are an option for missing excitation, not an intrinsic requirement
of implicit differential learning.

\subsection{A kinematic ambiguity in the periodic dynamic context}
The periodic plant used in the retained Fourier studies has
\begin{equation}
 \begin{aligned}
 f(z)&=\rho(\phi)r+\kappa(\phi)x-u=0,\\
 z&=(x,r,u,\phi)^{\mathsf T},\qquad
 r=\dot x,\quad\dot\phi=\omega\ne0.
 \end{aligned}
 \label{eq:tv_periodic_context}
\end{equation}
Its physical differential, written here for analysis only, is
\begin{equation}
 F(z)=\bigl(\kappa(\phi),\rho(\phi),-1,
                \rho'(\phi)r+\kappa'(\phi)x\bigr).
 \label{eq:tv_periodic_differential}
\end{equation}
An unrestricted direct rate field also admits the competing row
\begin{equation}
 H_{\rm kin}(z)=(\omega,0,0,-r),\qquad
 H_{\rm kin}(z)\dot z=\omega\dot x-r\dot\phi=0.
 \label{eq:tv_kinematic_competitor}
\end{equation}
This identity holds for every physical-time trajectory with these
kinematics wherever ordinary time derivatives exist, irrespective of the
input history. At held-input jumps, an ordinary $\dot u$ observation is
unavailable; the present differential calculation uses smooth histories.
The row has zero $r$ and $u$ entries and hence cannot complete the desired
$\dot r$ or $\dot u$ query.
It is nonzero and can be normalized, so excluding the zero field does not
exclude it. It is not itself a Jacobian on an open context region:
$\partial_r(H_{\rm kin})_\phi=-1$ whereas
$\partial_\phi(H_{\rm kin})_r=0$. A learning scheme that explicitly enforces
generic integrability could exclude this particular field; mere
normalization or the assumption that the physical target is a Jacobian
does not enforce that restriction on a learned field.

Now perturb a whole input history,
$u_\varepsilon(t)=u(t)+\varepsilon p(t)$, using the same clock and admissible
initial conditions. The plant induces $x_\varepsilon$ and
$r_\varepsilon=\partial_t x_\varepsilon$; these are not independently
assigned trial coordinates. At $\varepsilon=0$,
\begin{equation}
 \delta\phi=0,\qquad \delta r=\partial_t\delta x,\qquad
 \kappa\delta x+\rho\delta r-\delta u=0.
 \label{eq:tv_periodic_trial}
\end{equation}
Thus $F\delta z=0$, but
\begin{equation}
 H_{\rm kin}\delta z=\omega\delta x,
 \label{eq:tv_periodic_separation}
\end{equation}
which is nonzero wherever the input perturbation produces a nonzero state
variation. Such observations distinguish this competitor without supplying
the physical partial formulas. Excluding one competitor is not a proof of
complete identification; the combined admissible directions must still
meet the required span or lifted excitation condition.

Changing temporal input histories alone cannot eliminate the exact identity
(\ref{eq:tv_kinematic_competitor}) from an unrestricted context-dependent
candidate class. This does not rule out local temporal tracking under
additional shared-field assumptions. The distinction concerns the data
directions and the allowed field class, not the number of random seeds.

\subsection{Measurement and experimental consequences}
The exact variation in (\ref{eq:tv_family}) is a derivative with respect to
$\varepsilon$. A practical trial difference
\begin{equation}
 v_\varepsilon(t)=
       \frac{z_\varepsilon(t)-z_0(t)}{\varepsilon}
 \label{eq:tv_finite_trial}
\end{equation}
approximates it. Under twice differentiable dependence on $\varepsilon$,
$v_\varepsilon=\delta z+O(\varepsilon)$ and
$F(z_0)v_\varepsilon=O(\varepsilon)$. The linear example above is exact;
generic nonlinear trial differences are not exact tangent observations.
Reducing the perturbation also amplifies measurement-difference noise by
$1/|\varepsilon|$. Trial alignment, clock/initial-condition constraints,
perturbation size, sensor errors, and the resulting approximation bias must
therefore be documented. Finite temporal secants have their own curvature
and differentiation errors.

Real-time demonstrations can retain temporal observations whenever their
local or lifted excitation and tracking conditions are supported. If
additional null directions remain, possible remedies include improved
online excitation, additional admissible trajectory variations, or an
explicitly justified restriction on the learned field. None should be
silently substituted for the reported protocol. For the intended Fourier
Memory demonstration, partials must vary faster than the supplied inputs,
and this separation must be checked on the actual signals. Fourier Memory
represents fast partial variation;
it does not supply missing informative directions. Direct harmonic weights
parameterize independently learned partials, rather than a prescribed
physical residual whose coefficients are fitted and then differentiated.
Every demonstration must separately evaluate incompatibility, partial-field
recovery up to gauge, and readout using the same learned field.
